\begin{thebibliography}{99}

\bibitem{ref1} 祝树金, 朱捷尧. 数字普惠金融如何影响中国中小微企业的价值链攀升[J]. 上海大学学报(社会科学版), 2025, 42(4): 113-130.

\bibitem{ref2} 崔惠玉, 徐颖, 李秀华. 财政资金如何引导地方政府促进小微企业创新[J]. 财贸经济, 2026, 47(3). [页码待确认]

\bibitem{ref3} 王贤彬, 许婷君, 严祥起. 普惠小微信贷供给与小微企业进入——来自小微支行设立的证据[J]. 金融研究, 2025, 546(12): 95-113.

\bibitem{ref4} 吴勇民. 金融大模型赋能中小微企业供应链融资的双重机制与演化博弈研究[J]. 工业技术经济, 2026, 45(4): 16-29.

\bibitem{ref5} 汪勇, 冯心歌, 赵宸宇. 数字金融、信息不对称与小微企业融资——来自助贷机构的证据[J]. 数量经济技术经济研究, 2026, 43(1). [页码待确认]

\bibitem{ref6} 花青青, 王薪镐, 丁美玲. 基于SERVQUAL-IPA模型的金融机构服务质量评价与提升策略研究——以P银行J支行为例[J]. 知识经济, 2025(26): 65-69.

\bibitem{ref7} 周雷, 殷凯丽, 应皓恬, 等. 数字普惠金融服务小微企业融资研究——以全国首个小微企业数字征信实验区为例[J]. 西南金融, 2024(1): 54-68.

\bibitem{ref8} 韩莉, 宋路杰, 张杨林, 等. 金融科技如何助力小微企业融资——文献评析与展望[J]. 中国软科学, 2021(S01): 287-296.

\bibitem{ref9} 顾锋娟, 金德环, 曾守桢. 数字普惠金融促进小微企业融资的理论机理及省级证据[J]. 商业经济与管理, 2024, 387(1): 92-104.

\bibitem{ref10} 宋卓霖, 赵涵. 政策性担保能否破解小微企业融资难?——基于融担基金再担保合作的准自然实验[J]. 东方论坛, 2024(5): 85-101.

\bibitem{ref11} 田越. 邮储银行J支行对公业务客户满意度评价研究[D]. 石家庄: 河北地质大学, 2023.

\bibitem{ref12} 孙彦佩. ZY银行郑州分行网点服务质量提升策略研究[D]. 郑州: 河南财经政法大学, 2025.

\bibitem{ref13} 孙启蒙. 农行T支行基于精益六西格玛的零售业务服务质量改善研究[D]. 济南: 山东大学, 2023.

\bibitem{ref14} 唐家驰. 基于六西格玛理论的N银行Y分行网点服务管理优化研究[D]. 广州: 广东工业大学, 2025.

\bibitem{ref15} 张磊, 许坤, 张琳, 等. 政策扶持、金融科技与小微企业信贷融资[J]. 统计研究, 2023, 40(12): 50-61. DOI: 10.19343/j.cnki.11-1302/c.2023.12.005.

\bibitem{ref16} 黄玉婷. 基于精益六西格玛管理的M公司客户服务流程优化研究[D]. 深圳: 深圳大学, 2023.

\bibitem{ref17} 张小敏. A银行杭州分行小微企业线上贷款流程优化研究[D]. 上海: 华东师范大学, 2024. [学位授予单位待CNKI确认]

\bibitem{ref18} 黄兆锟. M银行P支行小微企业客户关系管理优化研究[D]. 南昌: 江西财经大学, 2025.

\bibitem{ref19} 李欣春. A银行小微企业普惠金融信贷业务营销策略优化研究[D]. 南京: 南京工业大学, 2025.

\bibitem{ref20} 金波, 牛华伟. 传统银行借贷、金融科技与中小企业融资[J]. 运筹与管理, 2024, 33(3): 169-176. DOI: 10.12005/orms.2024.0094.

\bibitem{ref21} Parasuraman A, Zeithaml V A, Berry L L. SERVQUAL: A multiple-item scale for measuring consumer perceptions of service quality[J]. Journal of Retailing, 1988, 64(1): 12-40.

\bibitem{ref22} Yao Y, Yang Z. Can digital finance boost SME innovation by easing financing constraints? Evidence from Chinese GEM-listed companies[J]. PLoS ONE, 2022, 17(3): e0264647.

\bibitem{ref23} Lu Z, Wu J, Li H, et al. Local bank, digital financial inclusion and SME financing constraints: empirical evidence from China[J]. Emerging Markets Finance and Trade, 2022, 58(6): 1712-1725. DOI: 10.1080/1540496X.2021.1923477.

\bibitem{ref24} Zouari G, Abdelhedi M. Customer satisfaction in the digital era: evidence from Islamic banking[J]. Journal of Innovation and Entrepreneurship, 2021, 10: 9. DOI: 10.1186/s13731-021-00151-x.

\bibitem{ref25} Kaur B, Kiran S, Grima S, Rupeika-Apoga R. Digital banking in Northern India: the risks on customer satisfaction[J]. Risks, 2021, 9(11): 209. DOI: 10.3390/risks9110209.

\bibitem{ref26} Nguyen D T, Pham V T, Tran D M, et al. Impact of service quality, customer satisfaction and switching costs on customer loyalty[J]. Journal of Asian Finance, Economics and Business, 2020, 7(8): 395-405.

\bibitem{ref27} Siddiqi K O. Interrelations between service quality attributes, customer satisfaction and customer loyalty in the retail banking sector in Bangladesh[J]. International Journal of Business and Management, 2011, 6(3): 12-36.

\bibitem{ref28} Kheng L L, Mahamad O, Ramayah T, et al. The impact of service quality on customer loyalty: a study of banks in Penang, Malaysia[J]. International Journal of Marketing Studies, 2010, 2(2): 57-66.

\bibitem{ref29} Raza S A, Umer A, Qureshi M A, et al. Internet banking service quality, e-customer satisfaction and loyalty: the modified e-SERVQUAL model[J]. The TQM Journal, 2020, 32(6): 1443-1466. DOI: 10.1108/TQM-02-2020-0019.

\bibitem{ref30} Nguyen T Q, Ngo L T T, Huynh N T, Quoc T L, Hoang L V. Assessing port service quality: an application of the extension fuzzy AHP and importance-performance analysis[J]. PLoS ONE, 2022, 17(2): e0264590. DOI: 10.1371/journal.pone.0264590.

\bibitem{ref31} Moslem S, Farooq D, Ghorbanzadeh O. A systematic review of analytic hierarchy process applications to solve transportation problems[J]. Information, 2023, 14(2): 73.

\bibitem{ref32} Adanigbo E, et al. Lean Six Sigma framework for reducing operational delays in customer support centers for fintech products[J]. International Journal for Multidisciplinary Research (IJFMR), 2021, 2(2): 160-167. DOI: 10.54660/.IJFMR.2021.2.2.160-167.

\bibitem{ref33} de Koning H, Does R J M M, Bisgaard S. Lean Six Sigma in financial services[J]. International Journal of Six Sigma and Competitive Advantage, 2008, 4(1): 1-17.

\bibitem{ref34} Guarraia P, Schwedel A. For banks in need: getting more from Lean Six Sigma[R]. Bain \& Company Brief, 2008.

\bibitem{ref35} Delahoz-Dominguez E, Mendoza A, Zuluaga R. A Six Sigma and DEA framework for quality assessment in banking services[J]. Administrative Sciences, 2024, 14(11): 295. DOI: 10.3390/admsci14110295.

\bibitem{ref36} Dumitrescu E, Hué S, Hurlin C, et al. Machine learning for credit scoring: improving logistic regression with non-linear decision-tree effects[J]. European Journal of Operational Research, 2022, 297(3): 1178-1192. DOI: 10.1016/j.ejor.2021.06.053.

\bibitem{ref37} Nugraha D P, Setiawan B, Nathan R J, Fekete-Farkas M. Fintech adoption drivers for innovation for SMEs in Indonesia[J]. Journal of Open Innovation: Technology, Market, and Complexity, 2022, 8(4): 208. DOI: 10.3390/joitmc8040208.

\bibitem{ref38} Zhang L, Chen J, Liu Z, Hao Z. Digital inclusive finance, financing constraints, and technological innovation of SMEs: differences in the effects of financial regulation and government subsidies[J]. Sustainability, 2023, 15(9): 7144. DOI: 10.3390/su15097144.

\bibitem{ref39} Martilla J A, James J C. Importance-performance analysis[J]. Journal of Marketing, 1977, 41(1): 77-79.

\bibitem{ref40} Yin R K. Case study research and applications: design and methods[M]. 6th ed. Thousand Oaks, CA: Sage Publications, 2018.

\bibitem{ref41} Stiglitz J E, Weiss A. Credit rationing in markets with imperfect information[J]. The American Economic Review, 1981, 71(3): 393-410.

\bibitem{ref42} Cronin J J, Taylor S A. Measuring service quality: a reexamination and extension[J]. Journal of Marketing, 1992, 56(3): 55-68.

\bibitem{ref43} Heskett J L, Jones T O, Loveman G W, et al. Putting the service-profit chain to work[J]. Harvard Business Review, 1994, 72(2): 164-174.

\bibitem{ref44} Davis F D. Perceived usefulness, perceived ease of use, and user acceptance of information technology[J]. MIS Quarterly, 1989, 13(3): 319-340.

\bibitem{ref45} Vashishth A, Lameijer B A, Chakraborty A, Antony J, Moormann J. Implementing Lean Six Sigma in financial services: the effect of motivations, selected methods and challenges on LSS program- and organizational performance[J]. International Journal of Quality \& Reliability Management, 2024, 41(2): 509-531. DOI: 10.1108/IJQRM-05-2022-0154.

\bibitem{ref46} 国务院. 推进普惠金融发展规划(2016—2020年)[EB/OL]. 2015.

\bibitem{ref47} 中国人民银行, 等. 关于进一步强化中小微企业金融服务的指导意见[EB/OL]. 2020.

\bibitem{ref48} Project Management Institute. PMBOK Guide: a guide to the project management body of knowledge[S]. 7th ed. Newtown Square, PA: PMI, 2021.

\bibitem{ref49} Zaheer A. Structure equation modeling (SEM) approach for evaluating and analyzing the effect of IT-based services in banking sector on customer service quality (SERVQUAL)[J]. International Journal of Advanced and Applied Sciences, 2023, 10(2): 147-155. DOI: 10.21833/ijaas.2023.02.018.

\bibitem{ref50} Agrawal V, Seth N, Dixit J K. A combined AHP-TOPSIS-DEMATEL approach for evaluating success factors of e-service quality: an experience from Indian banking industry[J]. Electronic Commerce Research, 2022, 22(3): 715-747. DOI: 10.1007/s10660-020-09430-3.

\bibitem{ref51} Liu J Y, Liu M H, Sun R Q. Quality management of financial shared service center in the digitalization context: a case study of HX Financial Shared Services Center[J]. Frontiers of Business Research in China, 2022, 16(2): 207-224. DOI: 10.3868/s070-007-022-0010-5.

\end{thebibliography}